\documentclass[11pt]{scrreprt}
\usepackage[sexy,noasy]{evan}
\input{preamble.sty}

% Cover and contents configuration
\setcounter{tocdepth}{1}
\setcounter{secnumdepth}{1}
\renewcommand{\contentsname}{Contents}

\begin{document}

% --------------------------------------------------
% Cover page (page 1)
% --------------------------------------------------
\begin{titlepage}
    \centering
    \vspace*{\fill}
    
    {\LARGE\bfseries CSE331 Proofs\par}

    \vspace{0.5cm}

    {\large Automata and Computability\par}

    \vspace{1.5cm}

    {\large\scshape Irfan Rahman\par}

    \vspace{0.7cm}

    {\small November 2026\par}

    \vspace*{\fill}
\end{titlepage}

% --------------------------------------------------
% Table of contents (page 2)
% --------------------------------------------------
\pagenumbering{roman}
\setcounter{page}{1}
\tableofcontents
\clearpage

% --------------------------------------------------
% Part I: Midterm
% --------------------------------------------------
\pagenumbering{arabic}
\setcounter{page}{1}
\chapter{Midterm}



% Union
\section{Closure Under Union}
\begin{theorem*}[Closure Under Union]
    The set of regular languages is closed under union.
\end{theorem*}

\begin{proof}
    Let $\mathcal{R}_L$ be the set of all regular languages.

    We need to show that
    \[
        \forall L_1, L_2 \in \mathcal{R}_L,
        \qquad L_1 \cup L_2 \in \mathcal{R}_L.
    \]

    Since $L_1$ and $L_2$ are regular, there exist two DFAs
    \[
        D_1 = (Q_1, \Sigma, \delta_1, q_1, F_1)
    \]
    and
    \[
        D_2 = (Q_2, \Sigma, \delta_2, q_2, F_2)
    \]
    that recognize $L_1$ and $L_2$, respectively.

    We construct a new DFA
    \[
        D = (Q, \Sigma, \delta, q_0, F)
    \]
    that recognizes $L_1 \cup L_2$, where:

    \begin{enumerate}
        \item $Q = Q_1 \times Q_2$.

        \item $\Sigma$ is the common input alphabet
        of $D_1$ and $D_2$.

        \item The transition function $\delta$ is defined as
        \[
            \delta((r_1,r_2),a)
            =
            (\delta_1(r_1,a),\delta_2(r_2,a))
        \]
        for all $r_1 \in Q_1$, $r_2 \in Q_2$,
        and $a \in \Sigma$.

        \item $q_0 = (q_1,q_2)$.

        \item The set of accepting states is
        \[
            F =
            \left\{
                (r_1,r_2)
                \mid
                r_1 \in F_1
                \text{ or }
                r_2 \in F_2
            \right\}.
        \]

        Equivalently,
        \[
            F =
            (F_1 \times Q_2)
            \cup
            (Q_1 \times F_2).
        \]
    \end{enumerate}

    The DFA $D$ simulates both $D_1$ and $D_2$
    simultaneously and accepts a string if
    either $D_1$ or $D_2$ accepts it.

    Therefore,
    \[
        L(D) = L_1 \cup L_2.
    \]

    Since $D$ is a DFA that recognizes
    $L_1 \cup L_2$, the language
    $L_1 \cup L_2$ is regular.

    Hence,
    \[
        \boxed{
            \mathcal{R}_L
            \text{ is closed under union.}
        }
    \]
\end{proof}

% Complement
\section{Closure Under Complement}
\begin{theorem*}[Closure Under Complement]
    The set of regular languages is closed under complement.
\end{theorem*}

\begin{proof}
    Let $\mathcal{R}_L$ be the set of all regular languages.

    We need to show that
    \[
        \forall L \in \mathcal{R}_L,
        \qquad \overline{L} \in \mathcal{R}_L.
    \]

    Since $L$ is regular, there exists a DFA
    \[
        M = (Q, \Sigma, \delta, q_0, F)
    \]
    that recognizes $L$.

    We construct a new DFA
    \[
        M' = (Q', \Sigma', \delta', q_0', F')
    \]
    that recognizes $\overline{L}$, where:

    \begin{enumerate}
        \item $Q' = Q$.

        \item $\Sigma' = \Sigma$.

        \item $\delta'(q,a) = \delta(q,a)$,
        for all $q \in Q$ and $a \in \Sigma$.

        \item $q_0' = q_0$.

        \item $F' = Q - F$.
    \end{enumerate}

    The DFA $M'$ has the same states, alphabet,
    transitions, and start state as $M$, but its
    accepting and non-accepting states are interchanged.

    Therefore,
    \[
        L(M') = \overline{L} = \Sigma^* - L.
    \]

    Since $M'$ is a DFA that recognizes
    $\overline{L}$, the language $\overline{L}$
    is regular.

    Hence,
    \[
        \boxed{
            \mathcal{R}_L
            \text{ is closed under complement.}
        }
    \]
\end{proof}

% Intersection
\section{Closure Under Intersection}
\begin{theorem*}[Closure Under Intersection]
    The set of regular languages is closed under intersection.
\end{theorem*}

\begin{proof}
    Let $\mathcal{R}_L$ be the set of all regular languages.

    We need to show that
    \[
        \forall L_1, L_2 \in \mathcal{R}_L,
        \qquad L_1 \cap L_2 \in \mathcal{R}_L.
    \]

    Since $L_1$ and $L_2$ are regular, there exist two DFAs
    \[
        D_1 = (Q_1, \Sigma, \delta_1, q_1, F_1)
    \]
    and
    \[
        D_2 = (Q_2, \Sigma, \delta_2, q_2, F_2)
    \]
    that recognize $L_1$ and $L_2$, respectively.

    We construct a new DFA
    \[
        D = (Q, \Sigma, \delta, q_0, F)
    \]
    that recognizes $L_1 \cap L_2$, where:

    \begin{enumerate}
        \item $Q = Q_1 \times Q_2$.

        \item $\Sigma$ is the common input alphabet of $D_1$ and $D_2$.

        \item The transition function $\delta$ is defined as
        \[
            \delta((r_1,r_2),a)
            =
            (\delta_1(r_1,a),\delta_2(r_2,a))
        \]
        for all $r_1 \in Q_1$, $r_2 \in Q_2$, and $a \in \Sigma$.

        \item $q_0 = (q_1,q_2)$.

        \item The set of accepting states is
        \[
            F =
            \left\{
                (r_1,r_2)
                \mid
                r_1 \in F_1 \text{ and } r_2 \in F_2
            \right\}.
        \]

        Equivalently,
        \[
            F = F_1 \times F_2.
        \]
    \end{enumerate}

    The DFA $D$ simulates both $D_1$ and $D_2$ simultaneously
    and accepts a string if both $D_1$ and $D_2$ accept it.

    Therefore,
    \[
        L(D) = L_1 \cap L_2.
    \]

    Since $D$ is a DFA that recognizes $L_1 \cap L_2$,
    the language $L_1 \cap L_2$ is regular.

    Hence,
    \[
        \boxed{
            \mathcal{R}_L
            \text{ is closed under intersection.}
        }
    \]
\end{proof}

% Difference
\section{Closure Under Difference}
\begin{theorem*}[Closure Under Difference]
    The set of regular languages is closed under difference.
\end{theorem*}

\begin{proof}
    Let $\mathcal{R}_L$ be the set of all regular languages.

    We need to show that
    \[
        \forall L_1, L_2 \in \mathcal{R}_L,
        \qquad L_1 - L_2 \in \mathcal{R}_L.
    \]

    Since $L_1$ and $L_2$ are regular, there exist two DFAs
    \[
        D_1 = (Q_1, \Sigma, \delta_1, q_1, F_1)
    \]
    and
    \[
        D_2 = (Q_2, \Sigma, \delta_2, q_2, F_2)
    \]
    that recognize $L_1$ and $L_2$, respectively.

    We construct a new DFA
    \[
        D = (Q, \Sigma, \delta, q_0, F)
    \]
    that recognizes $L_1 - L_2$, where:

    \begin{enumerate}
        \item $Q = Q_1 \times Q_2$.

        \item $\Sigma$ is the common input alphabet of $D_1$ and $D_2$.

        \item The transition function $\delta$ is defined as
        \[
            \delta((r_1,r_2),a)
            =
            (\delta_1(r_1,a),\delta_2(r_2,a))
        \]
        for all $r_1 \in Q_1$, $r_2 \in Q_2$, and $a \in \Sigma$.

        \item $q_0 = (q_1,q_2)$.

        \item The set of accepting states is
        \[
            F =
            \left\{
                (r_1,r_2)
                \mid
                r_1 \in F_1 \text{ and } r_2 \notin F_2
            \right\}.
        \]

        Equivalently,
        \[
            F = F_1 \times (Q_2 - F_2).
        \]
    \end{enumerate}

    The DFA $D$ simulates both $D_1$ and $D_2$ simultaneously
    and accepts a string if $D_1$ accepts it and $D_2$ rejects it.

    Therefore,
    \[
        L(D) = L_1 - L_2.
    \]

    Since $D$ is a DFA that recognizes $L_1 - L_2$,
    the language $L_1 - L_2$ is regular.

    Hence,
    \[
        \boxed{
            \mathcal{R}_L
            \text{ is closed under difference.}
        }
    \]
\end{proof}

% Concatenation
\section{Closure Under Concatenation}

\begin{theorem*}[Closure Under Concatenation]
    The set of regular languages is closed under concatenation.
\end{theorem*}

\begin{proof}
    Let $\mathcal{R}_L$ be the set of all regular languages.

    We need to show that
    \[
        \forall L_1, L_2 \in \mathcal{R}_L,
        \qquad L_1 \circ L_2 \in \mathcal{R}_L.
    \]

    Since $L_1$ and $L_2$ are regular, there exist two NFAs
    \[
        N_1 = (Q_1, \Sigma, \delta_1, q_1, F_1)
    \]
    and
    \[
        N_2 = (Q_2, \Sigma, \delta_2, q_2, F_2)
    \]
    that recognize $L_1$ and $L_2$, respectively.

    We construct a new NFA
    \[
        N = (Q, \Sigma, \delta, q_0, F)
    \]
    that recognizes $L_1 \circ L_2$, where:

    \begin{enumerate}
        \item $Q = Q_1 \cup Q_2$, assuming $Q_1$ and $Q_2$
        are disjoint.

        \item $\Sigma$ is the common input alphabet
        of $N_1$ and $N_2$.

        \item The transition function $\delta$ is defined as follows:

        For every $q \in Q$ and $a \in \Sigma_\varepsilon$,
        \[
            \delta(q,a) =
            \begin{cases}
                \delta_1(q,a),
                & q \in Q_1,\ q \notin F_1,\\[4pt]
                \delta_1(q,a),
                & q \in F_1,\ a \neq \varepsilon,\\[4pt]
                \delta_1(q,a) \cup \{q_2\},
                & q \in F_1,\ a = \varepsilon,\\[4pt]
                \delta_2(q,a),
                & q \in Q_2.
            \end{cases}
        \]

        \item $q_0 = q_1$.

        \item $F = F_2$.
    \end{enumerate}

    The NFA $N$ begins processing the input using $N_1$.
    Upon reaching an accepting state of $N_1$, it may
    move to the start state of $N_2$ through an
    $\varepsilon$-transition and process the remaining input.

    Thus, $N$ accepts a string if it can be divided
    into two parts such that $N_1$ accepts the first
    part and $N_2$ accepts the second part.

    Therefore,
    \[
        L(N) = L_1 \circ L_2.
    \]

    Since $N$ is an NFA that recognizes $L_1 \circ L_2$,
    the language $L_1 \circ L_2$ is regular.

    Hence,
    \[
        \boxed{
            \mathcal{R}_L
            \text{ is closed under concatenation.}
        }
    \]
\end{proof}    

% Kleene star
\section{Closure Under Kleene Star}
\begin{theorem*}[Closure Under Kleene Star]
    The set of regular languages is closed under Kleene star.
\end{theorem*}

\begin{proof}
    Let $\mathcal{R}_L$ be the set of all regular languages.

    We need to show that
    \[
        \forall L \in \mathcal{R}_L,
        \qquad L^* \in \mathcal{R}_L.
    \]

    Since $L$ is regular, there exists an NFA
    \[
        N_1 = (Q_1, \Sigma, \delta_1, q_1, F_1)
    \]
    that recognizes $L$.

    We construct a new NFA
    \[
        N = (Q, \Sigma, \delta, q_0, F)
    \]
    that recognizes $L^*$, where:

    \begin{enumerate}
        \item $Q = Q_1 \cup \{q_0\}$, where $q_0$ is a new state not in $Q_1$.

        \item $\Sigma$ is the same input alphabet as $N_1$.

        \item The transition function $\delta$ is defined as follows:

        For every $q \in Q$ and $a \in \Sigma_\varepsilon$,
        \[
            \delta(q,a)=
            \begin{cases}
                \delta_1(q,a), & q \in Q_1,\ a \in \Sigma,\\[4pt]
                \delta_1(q,\varepsilon) \cup \{q_1\},
                    & q \in F_1,\ a=\varepsilon,\\[4pt]
                \delta_1(q,\varepsilon),
                    & q \in Q_1 - F_1,\ a=\varepsilon,\\[4pt]
                \{q_1\}, & q=q_0,\ a=\varepsilon,\\[4pt]
                \varnothing, & q=q_0,\ a\in\Sigma.
            \end{cases}
        \]

        \item $q_0$ is the new start state.

        \item The set of accepting states is
        \[
            F = F_1 \cup \{q_0\}.
        \]
    \end{enumerate}

    The NFA $N$ introduces a new accepting start state $q_0$
    to accept the empty string $\varepsilon$. It also adds an
    $\varepsilon$-transition from $q_0$ to $q_1$ and from every
    accepting state of $N_1$ back to $q_1$.

    Thus, $N$ accepts the empty string or any string formed by
    concatenating zero or more strings accepted by $N_1$.

    Therefore,
    \[
        L(N)=L^*.
    \]

    Since $N$ is an NFA that recognizes $L^*$, the language
    $L^*$ is regular.

    Hence,
    \[
        \boxed{
            \mathcal{R}_L
            \text{ is closed under Kleene star.}
        }
    \]
\end{proof}



% --------------------------------------------------
% Part II: Final (add proof sections here later)
% --------------------------------------------------
\chapter{Final}

\noindent\emph{Final examination proofs will be added here.}

\end{document}
